\subsection{Temporal dynamics: the next falsification surface}
Most of the mechanism results above are single-turn, whereas real prompt-injection
attacks often unfold over multiple turns or tool-mediated states, so the next test is
temporal: whether the refusal order parameter $m_t = (h_t \cdot \dref)/\norm{h_t}$,
measured turn by turn at a fixed model and layer, collapses gradually, abruptly, or
recovers as an attack trajectory develops. A positive result would not establish a
human-like supervisory system; it would show that the residual projection whose
single-turn collapse is measured here also has structured dynamics across turns, while
a negative result would be equally informative, indicating that an attack can succeed
or fail by late decoding, tool policy, or context routing without a monotone internal
monitor collapse. The trajectory artifact is designed to be dual-use safe, public
outputs exclude harmful prompts, raw completions, injection payloads, direction
vectors, and checkpoints, and each trajectory names its model, layer, direction
source, attack family, and pilot or powered status, with controls comprising a benign
multi-turn task, a random-direction trajectory, and where possible a held-out attack
family such as AgentDojo indirect injection. The projection $m_t$ remains a
mechanism-tier quantity; describing its decline as symmetry-breaking or monitoring
suppression is analogy-tier language, and the binding constraint from the triangulation
applies, so no temporal collapse is read as behavioural capture until sufficiency and
dynamics are tested across models, layers, and attack families.
